% Source: https://artofproblemsolving.com/community/c3h1942235_2019_mathcounts_national_team_round_8
% Stable identity: aops-topic-1942235-q01
First $T=\dfrac{2019\cdot2020}{2}=2019\cdot1010=2\cdot3\cdot5\cdot101\cdot673$.

A sum of $k\ge1$ consecutive positive integers starting at $m\ge1$ equals $\dfrac{k(2m+k-1)}{2}$. So we need
\[
2T=k\cdot L,\qquad L=2m+k-1 .
\]
Here $L-k=2m-1$ is odd and positive, so $k$ and $L$ have opposite parity and $k<L$. Conversely, any factorisation $2T=k\cdot L$ with $k<L$ and $k,L$ of opposite parity gives $m=\dfrac{L-k+1}{2}$, a positive integer, and hence a representation.

Since $2T$ has exactly one factor $2$, in each factorisation into factors of opposite parity exactly one factor is odd, and it is an odd divisor $d$ of $T$; the other factor is $\dfrac{2T}{d}$. Each odd divisor $d$ gives one factorisation, and $k$ is the smaller of $d$ and $\dfrac{2T}{d}$. So the representations (with $k\ge1$) correspond exactly to the odd divisors of $T$. The odd part of $T$ is $3\cdot5\cdot101\cdot673$, which has $2^4=16$ divisors.

One of the $16$ representations has $k=1$ (the single number $T$, from $d=1$). Removing it leaves $\boxed{15}$ sums of two or more consecutive positive integers, including $1+2+\cdots+2019$ itself (from $d=2019$).
